% TOP_PROBLEM: {"id": 3159, "title": "Partitioning edge-connectivity", "category_id": 3, "category": {"id": 3, "name": "graph_theory", "display_name": "Graph Theory", "description": "Problems involving graphs, networks, and their properties.", "slug": "graph-theory", "order_index": 3, "created_at": "2026-07-31T15:26:25.671Z"}, "status": "open", "problem_number": "OPG-146", "set_id": 12}
% FIRST_POSTED: 2026-10-01
% ENTRY_KIND: statement_only
% UnsolvedMath ID 3159; source catalogue code OPG-146
% !TeX program = lualatex
% Definition and sources only; this file is not an LLM solution attempt.
\documentclass[11pt,a4paper]{article}
\usepackage[margin=25mm]{geometry}
\usepackage{fontspec}
\setmainfont{lmroman10-regular.otf}[BoldFont=lmroman10-bold.otf,
 ItalicFont=lmroman10-italic.otf,BoldItalicFont=lmroman10-bolditalic.otf]
\usepackage{amsmath,amssymb,mathtools}
\usepackage{unicode-math}
\setmathfont{latinmodern-math.otf}
\AtBeginDocument{\let\setminus\smallsetminus}
\usepackage{xurl,hyperref}
\hypersetup{colorlinks=true,urlcolor=blue}
\setcounter{secnumdepth}{0}
\setlength{\parindent}{0pt}
\setlength{\parskip}{5pt}
\setlength{\emergencystretch}{3em}
\begin{document}
\section*{Partitioning edge-connectivity}\label{3159}
{\small UnsolvedMath ID 3159 · OPG-146 · Graph Theory · OpenGarden}
\subsection{Definitions and mathematical statement}
Let $a,b\geq1$ be integers and let $G=(V,E)$ be a finite
undirected loopless multigraph on at least two vertices. Parallel
edges are counted separately. For a nonempty proper subset
$X\subset V$, write $\delta_G(X)$ for the set of edges having
exactly one endpoint in $X$. The graph is $k$-edge-connected
when $|\delta_G(X)|\geq k$ for every such $X$; equivalently,
deleting fewer than $k$ edges cannot disconnect it.

The question is whether the hypothesis
\[
 |\delta_G(X)|\geq a+b+2
 \qquad(\varnothing\neq X\subsetneq V)
\]
always permits a partition $E=A\mathbin{\dot\cup}B$ for which
$(V,A)$ is $a$-edge-connected and $(V,B)$ is $b$-edge-connected.
In terms of the same cuts, the desired simultaneous conditions are
\[
 |A\cap\delta_G(X)|\geq a,\qquad
 |B\cap\delta_G(X)|\geq b
 \qquad(\varnothing\neq X\subsetneq V).
\]
The two resulting graphs must be spanning: both retain every
vertex of $G$. Each original edge belongs to exactly one of
them, and the partition has to satisfy every cut constraint
at once. It may depend on $G,a,b$.

This asks for an additive connectivity threshold uniform over
all orders and edge multiplicities. Having enough edges in
each cut separately is only the input condition; the issue is
whether their assignments can be coordinated globally. The
special simple-graph case follows by restricting the allowed
multigraphs.
\subsection{Short English statement}
Can connectivity of at least a+b+2 be split into two spanning edge-disjoint graphs of edge connectivity at least a and b?
\subsection{Sources}
\begin{itemize}
\item[S1] ULAM.AI and the UnsolvedMath Contributors, supplied problems.json, record 3159 (OPG-146), OpenGarden collection. Catalogue identity and source transcription; CC BY 4.0.
\url{https://huggingface.co/datasets/ulamai/UnsolvedMath}.
\item[S2] Open Problem Garden, Partitioning edge-connectivity. Original problem and discussion; the mathematical formulation below is expanded from the supplied source record.
\url{http://www.openproblemgarden.org/op/partitioning_edge_connectivity}.
\end{itemize}
\end{document}
